\scanpage{48}
\deleted{\unsure{Supposons} $i$ \unsure{est égal à} $\unsure{0}$}

En tous cas, si $i=0$, la situation est décrite par le fait que l'on a fait une réduction du groupe structural de $\mathfrak S_3$ à $\mathbb Z/2$ \inserted{[$X''$ \inserted{\unsure{donc $X'$}} a une section !]}, et que \unsure{Tout} est déterminé par la donnée d'un élément de $H^1(S,\mathbb Z/2)$. \deleted{\unsure{Supposons}} \inserted{Et} si $i=1$, alors la situation est décrite en disant que l'on a une réduction du groupe structural à $\mathbb Z/3\subset\mathfrak S_3$ [\unsure{grâce} à une section de \inserted{la bisection} $X_1$] et la situation est décrite par la donnée d'un élément de $H^1(S,\mathbb Z/3)$.

\figure{\begin{tikzpicture}\draw (0,0)--(3,0);\end{tikzpicture}}

\deleted{\unsure{Donc}} En tous les cas encore, si $i$ est inversible, alors faisons opérer $\mathfrak S_3$ sur $P^1_S/(\mathbb Z/3)$ par passage au quotient, on trouve dans \deleted{\unsure{ce dernier}} \inserted{\unsure{ce fibré conique}} \deleted{\unsure{deux}} \uline{trois sections} distinctes, \inserted{(provenant de $X',\unsure{X''|X'''}$ et $s$)} que l'on peut choisir comme $\infty,0,1$, d'où une identification de ce dernier à $P^1_S$, et $X\longrightarrow P^1_S$ se factorise par
\diagram{\begin{tikzcd}[column sep=large,row sep=small]
X\arrow[r]&P^1_S\arrow[r]&P^1_S\\
&\infty,0,1\arrow[r]&\infty,0,i\\
&t\arrow[r,no head]&it^2
\end{tikzcd}}
; la bisection naturelle en question s'identifie alors à $\sqrt i$ ;
\mspage{47}
